\documentclass[acmsmall,nonacm]{acmart}
% enable page numbers
% \settopmatter{printfolios=true}
\settopmatter{printfolios=false}
%\documentclass[sigconf]{acmart}
%\settopmatter{printfolios=true,printccs=false,printacmref=false}
\input{../template/macro/pluverse-preamble.tex}
\newcommand{\lean}{\projectname{Lean4}}
\newcommand{\llvm}{\projectname{LLVM}}
\newcommand{\alive}{\projectname{Alive2}}
\newcommand{\smt}{\projectname{SMT}}
\newcommand{\vellvm}{\projectname{Vellvm}}
\newcommand{\leanmlir}{\projectname{LeanMLIR}}

\newcommand{\instcombine}{\textsc{InstCombine}\xspace}
\newcommand{\ir}{IR\xspace}
\newcommand{\poison}{\texttt{poison}\xspace}
\newcommand{\vundef}{\texttt{undef}\xspace}
\input{../template/macro/pluverse-preamble-end.tex}

\newif\ifInCopypaper
\InCopypaperfalse

%%
%% \BibTeX command to typeset BibTeX logo in the docs
\AtBeginDocument{%
  \providecommand\BibTeX{{%
    \normalfont B\kern-0.5em{\scshape i\kern-0.25em b}\kern-0.8em\TeX}}}

%% Rights management information.  This information is sent to you
%% when you complete the rights form.  These commands have SAMPLE
%% values in them; it is your responsibility as an author to replace
%% the commands and values with those provided to you when you
%% complete the rights form.
\copyrightyear{2026}
\acmYear{2026}
\setcopyright{cc}
\setcctype{by}
\acmConference[ASPLOS '26]{Proceedings of the 31st ACM International Conference on Architectural Support for Programming Languages and Operating Systems, Volume 2}{March 22--26, 2026}{Pittsburgh, PA, USA}
\acmBooktitle{Proceedings of the 31st ACM International Conference on Architectural Support for Programming Languages and Operating Systems, Volume 2 (ASPLOS '26), March 22--26, 2026, Pittsburgh, PA, USA}
\acmPrice{}
\acmDOI{10.1145/3779212.3790184}
\acmISBN{979-8-4007-2359-9/2026/03}


%%
%% Submission ID.
%% Use this when submitting an article to a sponsored event. You'll
%% receive a unique submission ID from the organizers
%% of the event, and this ID should be used as the parameter to this command.
%%\acmSubmissionID{123-A56-BU3}

%%
%% The majority of ACM publications use numbered citations and
%% references.  The command \citestyle{authoryear} switches to the
%% "author year" style.
%%
%% If you are preparing content for an event
%% sponsored by ACM SIGGRAPH, you must use the "author year" style of
%% citations and references.
%% Uncommenting
%% the next command will enable that style.
%%\citestyle{acmauthoryear}

%%
%% end of the preamble, start of the body of the document source.
\begin{document}

%%
%% The "title" command has an optional parameter,
%% allowing the author to define a "short title" to be used in page headers.
\title{Proposal: Formalization of LLVM Peephole Optimizations in Lean4}

%%
%% The "author" command and its associated commands are used to define
%% the authors and their affiliations.
%% Of note is the shared affiliation of the first two authors, and the
%% "authornote" and "authornotemark" commands
%% used to denote shared contribution to the research.
\author{Hongxu Xu}
\authornotemark[1]
\affiliation{%
	\institution{Cheriton School of Computer Science, University of Waterloo}
	\city{Waterloo}
	\country{Canada}}
\email{hongxu.xu@uwaterloo.ca}

%%
%% By default, the full list of authors will be used in the page
%% headers. Often, this list is too long, and will overlap
%% other information printed in the page headers. This command allows
%% the author to define a more concise list
%% of authors' names for this purpose.
%\renewcommand{\shortauthors}{Trovato and Tobin, et al.}

%%
%% The abstract is a short summary of the work to be presented in the
%% article.
\begin{abstract}
Formal methods have been increasingly applied to verify the correctness of compiler optimizations,
ensuring that transformations preserve program semantics.
This proposal outlines a project to formalize LLVM peephole optimizations in Lean4,
a powerful proof assistant.
By encoding the semantics of LLVM IR and the transformations performed by peephole optimizations,
we aim to create a framework for verifying the correctness of these optimizations.
Specifically, we will focus on a subset of LLVM IR which includes integer arithmetic and memory operations, and formalize a selection of common peephole optimizations.
Control flow is out of scope for this project, as we will focus on local tranformations within basic blocks.
This work will contribute to the reliability of compiler optimizations and provide insights into the design of future optimization passes.
\end{abstract}

%%
%% The code below is generated by the tool at http://dl.acm.org/ccs.cfm.
%% Please copy and paste the code instead of the example below.
%%
% \begin{CCSXML}
%     <ccs2012>
%     <concept>
%     <concept_id>10011007.10011006.10011041</concept_id>
%     <concept_desc>Software and its engineering~Compilers</concept_desc>
%     <concept_significance>500</concept_significance>
%     </concept>
%     </ccs2012>
% \end{CCSXML}

% \ccsdesc[500]{Software and its engineering~Compilers}

%%
%% Keywords. The author(s) should pick words that accurately describe
%% the work being presented. Separate the keywords with commas.
\keywords{Formal Methods, Compiler Optimization, LLVM, Lean4}


%%
%% This command processes the author and affiliation and title
%% information and builds the first part of the formatted document.
\maketitle

\section{Introduction}
\label{sec:intro}
\input{intro}

%%
%% The acknowledgments section is defined using the "acks" environment
%% (and NOT an unnumbered section). This ensures the proper
%% identification of the section in the article metadata, and the
%% consistent spelling of the heading.
% \begin{acks}
% \end{acks}

%%
%% The next two lines define the bibliography style to be used, and
%% the bibliography file.
\bibliographystyle{ACM-Reference-Format}
\balance
\bibliography{acmart}

\end{document}
\endinput
%%
%% End of file `sample-sigconf.tex'.